\chapter{MÉTODOS}
\section{Delineamento e responsabilidade}

Adotou-se revisão sistematizada, com busca documentada, critérios explícitos, seleção estruturada e síntese crítica descritiva, conforme a distinção de \textcite{grant2009}. O trabalho é de responsabilidade de um pesquisador, sem seleção, extração ou apreciação crítica em duplicata por revisores humanos independentes. Inteligência artificial apoiou a organização bibliográfica, as propostas de elegibilidade, a extração focal e a redação. A conferência autoral das decisões e estimativas desta versão ainda não está documentada.

O protocolo não foi pré-registrado. A emenda 0.7 documenta alterações após o conhecimento dos resultados: encerramento da recuperação, delimitação da busca de citações e alcance da extração e da apreciação narrativa. O histórico está no suplemento.

\section{Fontes, buscas e gerenciamento}

Consultaram-se Scopus e Web of Science Core Collection em 26/09/2026 e PubMed em 05/10/2026. Combinaram-se termos de material particulado, gravimetria, métodos automáticos, comparação e ar ambiente, sem filtros iniciais de ano, idioma ou tipo documental. PubMed combinou descritores \textit{Medical Subject Headings} (MeSH) e termos livres, sem restrição à MEDLINE. As estratégias executadas constam do suplemento e do repositório.

As bases produziram 233 registros, reduzidos a 150 pela deduplicação inicial. A busca de citações gerou 361 ocorrências e 287 novos registros: cinco sementes foram rastreadas retrospectiva e prospectivamente; uma sexta, apenas prospectivamente. Não se rastrearam todos os relatos selecionados. A base integrada reuniu 437 registros. A deduplicação complementar usou DOI e título exato normalizados; publicações de um mesmo estudo foram vinculadas quando identificadas.

Recuperaram-se textos em editoras, repositórios e documentos fornecidos pelo pesquisador. O encerramento documental ocorreu em 17/10/2026, sem reexecução das buscas ou reivindicação de exaustividade ou saturação.

\section{Elegibilidade e seleção}

Foram elegíveis resultados de ar ambiente externo que comparassem um método automático convencional, regulatório ou de pesquisa identificável com coleta em filtro e determinação gravimétrica, para \MPdez{} ou \MPdoiscinco{}. Exigiu-se o mesmo sítio representativo da mesma massa de ar, com períodos coincidentes ou explicitamente harmonizados. A designação ``referência'' não bastou para comprovar gravimetria.

Estudos mistos contribuíram por comparações externas, elegíveis e separáveis. Excluíram-se resultados de interiores, exposição pessoal ou laboratório; contrastes somente entre automáticos; faixas incompatíveis; sensores de baixo custo; e relatos sem contraste direto elegível. Admitiram-se relatórios, dissertações e trabalhos de congresso. O resumo de duas páginas de \textcite{rsl033} contribuiu qualitativamente; para o registro do congresso de Greene, utilizou-se a dissertação correspondente de \textcite{rsl003}.

Os registros foram avaliados por título e resumo e, quando disponíveis, pelo relato integral ou alternativo correspondente. Indisponibilidade do texto e identificação insuficiente foram registradas separadamente das exclusões por conteúdo.

\section{Extração e apreciação técnica}

A extração focal registrou contexto, métodos, pareamento, denominadores, resultados e limitações. Distinguiram-se resultados próprios de referências citadas, contrastes automático--gravimétrico de automático--automático e correlação $r$ de coeficiente de determinação $R^2$. Mantiveram-se a direção e a definição das medidas: razão entre médias, média de razões, inclinação e diferença absoluta permaneceram distintas. Os campos individuais constam do suplemento.

Os 78 estratos de relato e fração não enumeram todas as combinações de sítio, período e instrumento. Preservaram-se as agregações publicadas, sem assumir independência ou ponderar pelo número de observações. RSL-086 e RSL-107 foram vinculados à família de Pittsburgh; possíveis dependências adicionais constam da matriz.

A apreciação técnica narrativa examinou pareamento, gravimetria, operação e controle instrumental, denominadores, cobertura ambiental, estatística, transparência e participação de fabricantes. O piloto categórico permaneceu como histórico, sem classificação uniforme de risco ou escore global de qualidade para o conjunto ampliado.

\section{Síntese}

A síntese agrupou resultados por fração e tecnologia, distinguindo configurações e dados brutos ou corrigidos. O quadro apresenta contrastes focais; o suplemento, registros individuais. Não se calcularam efeito combinado, desvio médio global ou contagens por direção de efeito.

Correlações foram interpretadas como associação, sem demonstração isolada de concordância, conforme \textcite{bland1986}. Melhor desempenho nos dados usados para desenvolver uma correção não foi considerado validação externa. A gravimetria foi tratada como comparador operacional sujeito a incertezas e artefatos.
